\chapter{D'Artagnan Was To Meet an Old Acquaintance}

\lettrine{T}{here} is always something in a landing, if it be only from the smallest sea-boat—a trouble and a confusion which do not leave the mind the liberty of which it stands in need in order to study at the first glance the new locality presented to it. The moveable bridges, the agitated sailors, the noise of the water on the pebbles, the cries and importunities of those who wait upon the shores, are multiplied details of that sensation which is summed up in one single result—hesitation. It was not, then, till after standing several minutes on the shore that D'Artagnan saw upon the port, but more particularly in the interior of the isle, an immense number of workmen in motion. At his feet D'Artagnan recognized the five chalands laden with rough stone he had seen leave the port of Piriac. The smaller stones were transported to the shore by means of a chain formed by twenty-five or thirty peasants. The large stones were loaded on trollies which conveyed them in the same direction as the others, that is to say, towards the works, of which D'Artagnan could as yet appreciate neither the strength nor the extent. Everywhere was to be seen an activity equal to that which Telemachus observed on his landing at Salentum. D'Artagnan felt a strong inclination to penetrate into the interior; but he could not, under the penalty of exciting mistrust, exhibit too much curiosity. He advanced then little by little, scarcely going beyond the line formed by the fishermen on the beach, observing everything, saying nothing, and meeting all suspicion that might have been excited with a half-silly question or a polite bow. And yet, whilst his companions carried on their trade, giving or selling their fish to the workmen or the inhabitants of the city, D'Artagnan had gained by degrees, and, reassured by the little attention paid to him, he began to cast an intelligent and confident look upon the men and things that appeared before his eyes. And his very first glance fell on certain movements of earth about which the eye of a soldier could not be mistaken. At the two extremities of the port, in order that their fires should converge upon the great axis of the ellipse formed by the basin, in the first place, two batteries had been raised, evidently destined to receive flank pieces, for D'Artagnan saw the workmen finishing the platform and making ready the demi-circumference in wood upon which the wheels of the pieces might turn to embrace every direction over the épaulement. By the side of each of these batteries other workmen were strengthening gabions filled with earth, the lining of another battery. The latter had embrasures, and the overseer of the works called successively men who, with cords, tied the saucissons and cut the lozenges and right angles of turfs destined to retain the matting of the embrasures. By the activity displayed in these works, already so far advanced, they might be considered as finished: they were not yet furnished with their cannons, but the platforms had their gîtes and their madriers all prepared; the earth, beaten carefully, was consolidated; and supposing the artillery to be on the island, in less than two or three days the port might be completely armed. That which astonished D'Artagnan, when he turned his eyes from the coast batteries to the fortifications of the city, was to see that Belle-Isle was defended by an entirely new system, of which he had often heard the Comte de la Fère speak as a wonderful advance, but of which he had as yet never seen the application. These fortifications belonged neither to the Dutch method of Marollais, nor to the French method of the Chevalier Antoine de Ville, but to the system of Manesson Mallet, a skilful engineer, who about six or eight years previously had quitted the service of Portugal to enter that of France. The works had this peculiarity, that instead of rising above the earth, as did the ancient ramparts destined to defend a city from escalades, they, on the contrary, sank into it; and what created the height of the walls was the depth of the ditches. It did not take long to make D'Artagnan perceive the superiority of such a system, which gives no advantage to cannon. Besides, as the fosses were lower than, or on a level with, the sea, these fosses could be instantly inundated by means of subterranean sluices. Otherwise, the works were almost complete, and a group of workmen, receiving orders from a man who appeared to be conductor of the works, were occupied in placing the last stones. A bridge of planks thrown over the fosses for the greater convenience of the manœvers connected with the barrows, joined the interior to the exterior. With an air of simple curiosity D'Artagnan asked if he might be permitted to cross the bridge, and he was told that no order prevented it. Consequently he crossed the bridge, and advanced towards the group.

This group was superintended by the man whom D'Artagnan had already remarked, and who appeared to be the engineer-in-chief. A plan was lying open before him upon a large stone forming a table, and at some paces from him a crane was in action. This engineer, who by his evident importance first attracted the attention of D'Artagnan, wore a \binderyfrench{justaucorps}, which, from its sumptuousness, was scarcely in harmony with the work he was employed in, that rather necessitated the costume of a master-mason than of a noble. He was a man of immense stature and great square shoulders, and wore a hat covered with feathers. He gesticulated in the most majestic manner, and appeared, for D'Artagnan only saw his back, to be scolding the workmen for their idleness and want of strength.

D'Artagnan continued to draw nearer. At that moment the man with the feathers ceased to gesticulate, and, with his hands placed upon his knees, was following, half-bent, the effort of six workmen to raise a block of hewn stone to the top of a piece of timber destined to support that stone, so that the cord of the crane might be passed under it. The six men, all on one side of the stone, united their efforts to raise it to eight or ten inches from the ground, sweating and blowing, whilst a seventh got ready for when there should be daylight enough beneath it to slide in the roller that was to support it. But the stone had already twice escaped from their hands before gaining a sufficient height for the roller to be introduced. There can be no doubt that every time the stone escaped them, they bounded quickly backwards, to keep their feet from being crushed by the refalling stone. Every time, the stone, abandoned by them, sunk deeper into the damp earth, which rendered the operation more and more difficult. A third effort was followed by no better success, but with progressive discouragement. And yet, when the six men were bent towards the stone, the man with the feathers had himself, with a powerful voice, given the word of command, <Ferme!> which regulates manœvers of strength. Then he drew himself up.

<Oh! oh!> said he, <what is this all about? Have I to do with men of straw? \swear{Corne de bœuf!} stand on one side, and you shall see how this is to be done.>

<\swear{Peste!}> said D'Artagnan, <will he pretend to raise that rock? that would be a sight worth looking at.>

The workmen, as commanded by the engineer, drew back with their ears down, and shaking their heads, with the exception of the one who held the plank, who prepared to perform the office. The man with the feathers went up to the stone, stooped, slipped his hands under the face lying upon the ground, stiffened his Herculean muscles, and without a strain, with a slow motion, like that of a machine, lifted the end of the rock a foot from the ground. The workman who held the plank profited by the space thus given him, and slipped the roller under the stone.

<That's the way,> said the giant, not letting the rock fall again, but placing it upon its support.

<\swear{Mordioux!}> cried D'Artagnan, <I know but one man capable of such a feat of strength.>

<\swear{Hein!}> cried the colossus, turning round.

<Porthos!> murmured D'Artagnan, seized with stupor, <Porthos at Belle-Isle!>

On his part, the man with the feathers fixed his eyes upon the disguised lieutenant, and, in spite of his metamorphosis, recognized him. <D'Artagnan!> cried he; and the colour mounted to his face. <Hush!> said he to D'Artagnan.

<Hush!> in his turn, said the musketeer. In fact, if Porthos had just been discovered by D'Artagnan, D'Artagnan had just been discovered by Porthos. The interest of the particular secret of each struck them both at the same instant. Nevertheless the first movement of the two men was to throw their arms around each other. What they wished to conceal from the bystanders, was not their friendship, but their names. But, after the embrace, came reflection.

<What the devil brings Porthos to Belle-Isle, lifting stones?> said D'Artagnan; only D'Artagnan uttered that question in a low voice. Less strong in diplomacy than his friend, Porthos thought aloud.

<How the devil did you come to Belle-Isle?> asked he of D'Artagnan; <and what do you want to do here?> It was necessary to reply without hesitation. To hesitate in answer to Porthos would have been a check, for which the self-love of D'Artagnan would never have consoled itself.

<\swear{Pardieu!} my friend, I am at Belle-Isle because you are here.>

<Ah, bah!> said Porthos, visibly stupefied with the argument and seeking to account for it to himself, with the felicity of deduction we know to be particular to him.

<Without doubt,> continued D'Artagnan, unwilling to give his friend time to recollect himself, <I have been to see you at Pierrefonds.>

<Indeed!>

<Yes.>

<And you did not find me there?>

<No, but I found Mouston.>

<Is he well?>

<\swear{Peste!}>

<Well, but Mouston did not tell you I was here.>

<Why should he not? Have I, perchance, deserved to lose his confidence?>

<No; but he did not know it.>

<Well; that is a reason at least that does not offend my self-love.>

<Then how did you manage to find me?>

<My dear friend, a great noble like you always leaves traced behind him on his passage; and I should think but poorly of myself, if I were not sharp enough to follow the traces of my friends.> This explanation, flattering as it was, did not entirely satisfy Porthos.

<But I left no traces behind me, for I came here disguised,> said Porthos.

<Ah! You came disguised did you?> said D'Artagnan.

<Yes.>

<And how?>

<As a miller.>

<And do you think a great noble, like you, Porthos, can affect common manners so as to deceive people?>

<Well, I swear to you my friend, that I played my part so well that everybody was deceived.>

<Indeed! so well, that I have not discovered and joined you?>

<Yes; but how did you discover and join me?>

<Stop a bit. I was going to tell you how. Do you imagine Mouston\longdash>

<Ah! it was that fellow, Mouston,> said Porthos, gathering up those two triumphant arches which served him for eyebrows.

<But stop, I tell you—it was no fault of Mouston's because he was ignorant of where you were.>

<I know he was; and that is why I am in such haste to understand\longdash>

<Oh! how impatient you are, Porthos.>

<When I do not comprehend, I am terrible.>

<Well, you will understand. Aramis wrote to you at Pierrefonds, did he not?>

<Yes.>

<And he told you to come before the equinox.>

<That is true.>

<Well! that is it,> said D'Artagnan, hoping that this reason would mystify Porthos. Porthos appeared to give himself up to a violent mental labour.

<Yes, yes,> said he, <I understand. As Aramis told me to come before the equinox, you have understood that that was to join him. You then inquired where Aramis was, saying to yourself, <Where Aramis is, there Porthos will be.> You have learnt that Aramis was in Bretagne, and you said to yourself, <Porthos is in Bretagne.>>

<Exactly. In good truth, Porthos, I cannot tell why you have not turned conjuror. So you understand that, arriving at Roche-Bernard, I heard of the splendid fortifications going on at Belle-Isle. The account raised my curiosity, I embarked in a fishing boat, without dreaming that you were here: I came, and I saw a monstrous fine fellow lifting a stone Ajax could not have stirred. I cried out, <Nobody but the Baron de Bracieux could have performed such a feat of strength.> You heard me, you turned round, you recognized me, we embraced; and, \swear{ma foi!} if you like, my dear friend, we will embrace again.>

<Ah! now all is explained,> said Porthos; and he embraced D'Artagnan with so much friendship as to deprive the musketeer of his breath for five minutes.

<Why, you are stronger than ever,> said D'Artagnan, <and still, happily, in your arms.> Porthos saluted D'Artagnan with a gracious smile. During the five minutes D'Artagnan was recovering his breath, he reflected that he had a very difficult part to play. It was necessary that he always should question and never reply. By the time his respiration returned, he had fixed his plans for the campaign. 